% BACKEND - diseño de 27 casos (CP-001..CP-027)
% GUÍA EQUIPO: Responsable asignado: Saimol Jimenez (Backend).
\paragraph{Módulo Backend y reglas de negocio (27 casos: CP-001 a CP-027).}

\begin{disenio}
\caso{CP-001}{\textbf{Compra ADMIN para cliente con destinatario de regalo.} NR-01, NR-02; NP-01 (C01.1, Alta)}{Token ADMIN; CLIENT-A con contactPhone distinto al de giftRecipient; ítem en stock; pago CASH; Idempotency-Key UUIDv4.}{1) Enviar POST /\allowbreak{}api/\allowbreak{}v1/\allowbreak{}orders/\allowbreak{}admin con buyer REGISTERED\_CLIENT y giftRecipient. 2) Consultar GET /\allowbreak{}api/\allowbreak{}v1/\allowbreak{}orders/\allowbreak{}:\allowbreak{}id y /\allowbreak{}receipt (ADMIN). 3) Consultar GET /\allowbreak{}api/\allowbreak{}v1/\allowbreak{}shop/\allowbreak{}orders/\allowbreak{}:\allowbreak{}id (CLIENT-A).}{HTTP 201; pedido conserva comprador y destinatario separados; ADMIN lee pedido/\allowbreak{}comprobante y CLIENT-A consulta su compra.}{Titularidad asociada a CLIENT-A, datos de giftRecipient independientes, stock descontado y comprobante comercial coherente.}
\caso{CP-002}{\textbf{Rechazo compra ADMIN cliente sin teléfono de comprador.} NR-01, NR-02; NP-01 (C01.2, Alta)}{Token ADMIN; payload POST /\allowbreak{}api/\allowbreak{}v1/\allowbreak{}orders/\allowbreak{}admin con buyer.\allowbreak{}type=\allowbreak{}REGISTERED\_CLIENT omitiendo contactPhone; Idempotency-Key UUIDv4.}{1) Registrar conteo inicial de pedidos, comprobantes y stock del producto. 2) Enviar POST /\allowbreak{}api/\allowbreak{}v1/\allowbreak{}orders/\allowbreak{}admin sin contactPhone del comprador. 3) Consultar conteos y stock posterior.}{HTTP 400 de validación rechazando la creación del pedido por falta de contactPhone obligatorio del comprador.}{Sin pedido creado, sin PurchaseReceipt ni descuento de stock; conteos inicial y final idénticos.}
\caso{CP-003}{\textbf{Compra ADMIN invitado y lectura de comprobante.} NR-01, NR-02, NR-06, NR-09; NP-02 (C02.1, Alta)}{Token ADMIN; buyer.\allowbreak{}type=\allowbreak{}GUEST con name y contactPhone válidos; giftRecipient opcional; pago CASH; Idempotency-Key UUIDv4.}{1) Enviar POST /\allowbreak{}api/\allowbreak{}v1/\allowbreak{}orders/\allowbreak{}admin como invitado. 2) Invocar GET /\allowbreak{}api/\allowbreak{}v1/\allowbreak{}orders/\allowbreak{}:\allowbreak{}id/\allowbreak{}receipt con token ADMIN. 3) Verificar estructura imprimible del comprobante.}{HTTP 201 en creación; GET /\allowbreak{}receipt retorna 200 con número, ítems, importes en centavos, comprador invitado y aviso no fiscal.}{Comprobante corresponde exactamente a la venta y contiene los campos necesarios para impresión ADMIN sin exigir cuenta de usuario.}
\caso{CP-004}{\textbf{Rechazo compra ADMIN invitado sin teléfono obligatorio.} NR-01, NR-02, NR-06, NR-09; NP-02 (C02.2, Alta)}{Token ADMIN; buyer.\allowbreak{}type=\allowbreak{}GUEST con name válido y contactPhone omitido o vacío; Idempotency-Key UUIDv4.}{1) Verificar stock y total de pedidos/\allowbreak{}usuarios antes del envío. 2) Enviar POST /\allowbreak{}api/\allowbreak{}v1/\allowbreak{}orders/\allowbreak{}admin sin contactPhone del invitado. 3) Reconsultar stock y tablas.}{HTTP 400 rechazando el pedido de invitado sin teléfono obligatorio del comprador.}{No se persiste pedido ni comprobante, no se crea usuario ni se descuenta stock del producto.}
\caso{CP-005}{\textbf{Compra invitado con teléfono de cliente sin vínculo.} NR-01, NR-02, NR-06, NR-09; NP-02 (C02.3, Alta)}{Token ADMIN; buyer.\allowbreak{}type=\allowbreak{}GUEST cuyo contactPhone coincide exactamente con el teléfono de CLIENT-A; Idempotency-Key UUIDv4.}{1) Crear venta ADMIN de invitado con teléfono de CLIENT-A. 2) Consultar GET /\allowbreak{}api/\allowbreak{}v1/\allowbreak{}shop/\allowbreak{}orders y /\allowbreak{}shop/\allowbreak{}orders/\allowbreak{}:\allowbreak{}id con JWT de CLIENT-A. 3) Verificar cuentas en BD.}{Venta creada con userId=\allowbreak{}null; CLIENT-A no ve el pedido en su listado y su consulta directa devuelve 404/\allowbreak{}403.}{No se crea vínculo por coincidencia telefónica ni cuenta nueva; aislamiento de acceso preservado.}
\caso{CP-006}{\textbf{Pago externo declarado ADMIN en efectivo.} NR-04, NR-06; NP-02, NP-16 (C03.1, Alta)}{Token ADMIN; pedido válido con payment=\allowbreak{}\{method:\allowbreak{}'CASH', note:\allowbreak{}'Pago en caja'\}; Idempotency-Key UUIDv4; sin datos de tarjeta.}{1) Enviar POST /\allowbreak{}api/\allowbreak{}v1/\allowbreak{}orders/\allowbreak{}admin con payment.\allowbreak{}method=\allowbreak{}CASH. 2) Consultar respuesta, GET /\allowbreak{}api/\allowbreak{}v1/\allowbreak{}orders/\allowbreak{}:\allowbreak{}id/\allowbreak{}receipt y registro de pago declarado.}{HTTP 201 con kind=\allowbreak{}ADMIN\_DECLARED\_EXTERNAL, status=\allowbreak{}DECLARED, method=\allowbreak{}CASH, declaredById y declaredAt; sin débito simulado.}{Declaración y comprobante coherentes en CASH; cero campos bancarios o cobros de tarjeta simulada.}
\caso{CP-007}{\textbf{Pago externo declarado ADMIN por transferencia.} NR-04, NR-06; NP-02, NP-16 (C03.2, Alta)}{Token ADMIN; pedido válido con payment=\allowbreak{}\{method:\allowbreak{}'TRANSFER'\} sin referencia bancaria ni comprobante bancario; Idempotency-Key UUIDv4.}{1) Enviar POST /\allowbreak{}api/\allowbreak{}v1/\allowbreak{}orders/\allowbreak{}admin con payment.\allowbreak{}method=\allowbreak{}TRANSFER sin referencia bancaria. 2) Consultar GET /\allowbreak{}api/\allowbreak{}v1/\allowbreak{}orders/\allowbreak{}:\allowbreak{}id/\allowbreak{}receipt.}{HTTP 201 sin exigir referencia bancaria; status=\allowbreak{}DECLARED, method=\allowbreak{}TRANSFER, declaredById y declaredAt; comprobante sin verificación bancaria.}{Se persiste transferencia declarada por ADMIN sin presentar validación bancaria automática ni datos bancarios inventados.}
\caso{CP-008}{\textbf{Cambio de tasa IVA tras cotizar (409 SHOP\_SETTINGS\_CHANGED).} NR-05, NR-06; NP-01, NP-02, NP-03 (C04.1, Alta)}{Token CLIENT-A; carrito válido; rateBps inicial 1500 (15\%); PATCH admin/\allowbreak{}shop-settings cambia rateBps a 1200; Idempotency-Key UUIDv4.}{1) Ejecutar POST /\allowbreak{}api/\allowbreak{}v1/\allowbreak{}shop/\allowbreak{}checkout/\allowbreak{}quote como CLIENT-A. 2) Cambiar tax.\allowbreak{}rateBps vía PATCH /\allowbreak{}api/\allowbreak{}v1/\allowbreak{}admin/\allowbreak{}shop-settings. 3) Enviar POST /\allowbreak{}api/\allowbreak{}v1/\allowbreak{}shop/\allowbreak{}checkout con el quoteToken previo.}{POST /\allowbreak{}api/\allowbreak{}v1/\allowbreak{}shop/\allowbreak{}checkout responde HTTP 409 con error.\allowbreak{}code=\allowbreak{}SHOP\_SETTINGS\_CHANGED rechazando condiciones obsoletas.}{Sin pedido creado, sin PurchaseReceipt, sin descuento de stock ni débito simulado en la cuenta del cliente.}
\caso{CP-009}{\textbf{Cambio de precio tras cotizar e inmutabilidad 15\% IVA.} NR-05, NR-06; NP-01, NP-02, NP-03 (C04.2, Alta)}{Ítems por 7300 centavos, IVA 15\% (1095), total 8395; cambio de precio tras cotizar y cambio de tasa/\allowbreak{}precio post-compra.}{1) Cotizar, alterar precio y confirmar \ensuremath{\rightarrow} 409 PRICE\_CHANGED. 2) Recotizar subtotal 7300 y confirmar checkout (201). 3) Cambiar tasa/\allowbreak{}precio y releer GET /\allowbreak{}api/\allowbreak{}v1/\allowbreak{}shop/\allowbreak{}orders/\allowbreak{}:\allowbreak{}id/\allowbreak{}receipt.}{Primer intento retorna 409 PRICE\_CHANGED sin efectos; compra válida fija subtotal 7300, IVA 1095, total 8395 y aviso no fiscal.}{Relecturas del comprobante tras cambios posteriores muestran importes idénticos (7300/\allowbreak{}1095/\allowbreak{}8395) y sin etiqueta de factura fiscal.}
\caso{CP-010}{\textbf{Corrección autorizada de contacto ADMIN y auditoría.} NR-03, NR-07; NP-04 (C07.1, Alta)}{Token ADMIN; orderId existente; payload \{contactRevision:\allowbreak{}0, buyerPhone:\allowbreak{}'+593990001112', recipientPhone:\allowbreak{}'+593990001113'\}.}{1) Enviar PATCH /\allowbreak{}api/\allowbreak{}v1/\allowbreak{}orders/\allowbreak{}:\allowbreak{}id/\allowbreak{}contact con contactRevision=\allowbreak{}0. 2) Consultar GET /\allowbreak{}api/\allowbreak{}v1/\allowbreak{}orders/\allowbreak{}:\allowbreak{}id/\allowbreak{}contact-history.}{HTTP 200 actualiza contacto y eleva contactRevision a 1; contact-history lista cambios con field, previousValue, nextValue, actorId y createdAt.}{Corrección operativa aplicada y auditoría append-only trazable al actor ADMIN con valores anterior y nuevo.}
\caso{CP-011}{\textbf{Aislamiento de corrección de contacto en recibo y pedido.} NR-03, NR-07; NP-04 (C07.2, Alta)}{Token ADMIN; orderId P1 (contactRevision=\allowbreak{}0) con nuevo recipientPhone; orderId P2 de control sin cambios.}{1) Capturar snapshot de GET /\allowbreak{}orders/\allowbreak{}:\allowbreak{}id1/\allowbreak{}receipt y GET /\allowbreak{}orders/\allowbreak{}:\allowbreak{}id2. 2) Ejecutar PATCH /\allowbreak{}api/\allowbreak{}v1/\allowbreak{}orders/\allowbreak{}:\allowbreak{}id1/\allowbreak{}contact. 3) Reconsultar comprobante de P1 y datos de P2.}{P1 refleja contacto corregido; GET /\allowbreak{}orders/\allowbreak{}:\allowbreak{}id1/\allowbreak{}receipt mantiene snapshot original y P2 conserva sus datos y contactRevision sin cambio.}{Comparación antes/\allowbreak{}después demuestra aislamiento total: comprobante congelado y segundo pedido intactos.}
\caso{CP-012}{\textbf{Rechazo de corrección de contacto por rol no autorizado.} NR-03, NR-07; NP-04 (C07.3, Alta)}{Tokens JWT de CLIENT-A y DRIVER-A; orderId válido con contactRevision actual; payload válido de corrección de contacto.}{1) Enviar PATCH /\allowbreak{}api/\allowbreak{}v1/\allowbreak{}orders/\allowbreak{}:\allowbreak{}id/\allowbreak{}contact con token CLIENT-A. 2) Repetir con token DRIVER-A. 3) Consultar contacto y contact-history como ADMIN.}{Ambas peticiones retornan HTTP 403 Forbidden sin ejecutar la corrección.}{Contacto, contactRevision e historial de auditoría permanecen idénticos tras los dos intentos rechazados.}
\caso{CP-013}{\textbf{Horario ADMIN límite inicial 07:\allowbreak{}00 America/\allowbreak{}Guayaquil.} NR-14; NP-06 (C10.5, Alta)}{Token ADMIN; pedido ADMIN válido con scheduledDeliveryDate en fecha futura a las 07:\allowbreak{}00:\allowbreak{}00-05:\allowbreak{}00; Idempotency-Key UUIDv4.}{1) Verificar deliverySchedule 07:\allowbreak{}00-17:\allowbreak{}00 en shop-settings. 2) Enviar POST /\allowbreak{}api/\allowbreak{}v1/\allowbreak{}orders/\allowbreak{}admin con hora 07:\allowbreak{}00-05:\allowbreak{}00. 3) Consultar pedido y comprobante.}{HTTP 201; pedido creado en estado PENDING con hora local 07:\allowbreak{}00 en America/\allowbreak{}Guayaquil y comprobante emitido.}{Límite inicial inclusivo 07:\allowbreak{}00 aceptado y stock descontado correctamente.}
\caso{CP-014}{\textbf{Horario ADMIN límite final 17:\allowbreak{}00 America/\allowbreak{}Guayaquil.} NR-14; NP-06 (C10.6, Alta)}{Token ADMIN; pedido ADMIN válido con scheduledDeliveryDate en fecha futura a las 17:\allowbreak{}00:\allowbreak{}00-05:\allowbreak{}00; Idempotency-Key UUIDv4.}{1) Verificar deliverySchedule 07:\allowbreak{}00-17:\allowbreak{}00. 2) Enviar POST /\allowbreak{}api/\allowbreak{}v1/\allowbreak{}orders/\allowbreak{}admin con hora 17:\allowbreak{}00-05:\allowbreak{}00. 3) Consultar pedido y comprobante.}{HTTP 201; pedido creado en estado PENDING con hora local 17:\allowbreak{}00 en America/\allowbreak{}Guayaquil y comprobante generado.}{Límite final inclusivo 17:\allowbreak{}00 aceptado con reserva de stock y datos consistentes.}
\caso{CP-015}{\textbf{Rechazo horario ADMIN inválido 06:\allowbreak{}59 (422).} NR-14; NP-06 (C10.7, Alta)}{Token ADMIN; pedido ADMIN con scheduledDeliveryDate en fecha futura a las 06:\allowbreak{}59:\allowbreak{}00-05:\allowbreak{}00; Idempotency-Key UUIDv4.}{1) Registrar conteo de pedidos, comprobantes y stock. 2) Enviar POST /\allowbreak{}api/\allowbreak{}v1/\allowbreak{}orders/\allowbreak{}admin con hora 06:\allowbreak{}59-05:\allowbreak{}00. 3) Reconsultar conteos y stock.}{HTTP 422 con error.\allowbreak{}code=\allowbreak{}DELIVERY\_OUTSIDE\_HOURS rechazando el horario previo a 07:\allowbreak{}00.}{Sin pedido ni comprobante creado y stock del producto sin variación.}
\caso{CP-016}{\textbf{Rechazo horario ADMIN inválido 17:\allowbreak{}01 (422).} NR-14; NP-06 (C10.8, Alta)}{Token ADMIN; pedido ADMIN con scheduledDeliveryDate en fecha futura a las 17:\allowbreak{}01:\allowbreak{}00-05:\allowbreak{}00; Idempotency-Key UUIDv4.}{1) Registrar conteo de pedidos, comprobantes y stock. 2) Enviar POST /\allowbreak{}api/\allowbreak{}v1/\allowbreak{}orders/\allowbreak{}admin con hora 17:\allowbreak{}01-05:\allowbreak{}00. 3) Reconsultar conteos y stock.}{HTTP 422 con error.\allowbreak{}code=\allowbreak{}DELIVERY\_OUTSIDE\_HOURS rechazando el horario posterior a 17:\allowbreak{}00.}{Sin efectos parciales: conteos de pedidos, comprobantes y stock idénticos antes y después.}
\caso{CP-017}{\textbf{Agrupación temporal por extremos (10:\allowbreak{}00, 10:\allowbreak{}25, 10:\allowbreak{}50).} NR-18, NR-22; NP-08 (C13.1, Media)}{Token ADMIN; 3 pedidos PENDING elegibles hoy dentro de cobertura con horas 10:\allowbreak{}00, 10:\allowbreak{}25 y 10:\allowbreak{}50 (-05:\allowbreak{}00).}{1) Confirmar routePlanning.\allowbreak{}timeWindow\allowbreak{}Margin\allowbreak{}Minutes=\allowbreak{}30. 2) Invocar GET /\allowbreak{}api/\allowbreak{}v1/\allowbreak{}route-proposals?date=\allowbreak{}HOY como ADMIN. 3) Inspeccionar grupos y paradas.}{HTTP 200 con 2 propuestas separadas: [10:\allowbreak{}00, 10:\allowbreak{}25] y [10:\allowbreak{}50], porque 10:\allowbreak{}50 - 10:\allowbreak{}00 = 50 \ensuremath{>} 30 min (sin encadenamiento).}{Agrupación cumple max - min \ensuremath{\leq} 30 por extremos y orden cronológico de paradas sin persistir rutas.}
\caso{CP-018}{\textbf{Agrupación por extremos con separación de 29 minutos.} NR-18, NR-22; NP-08 (C13.2, Media)}{Token ADMIN; 2 pedidos PENDING elegibles hoy con horas 10:\allowbreak{}00 y 10:\allowbreak{}29 (-05:\allowbreak{}00), diferencia 29 minutos.}{1) Verificar margen=\allowbreak{}30. 2) Ejecutar GET /\allowbreak{}api/\allowbreak{}v1/\allowbreak{}route-proposals?date=\allowbreak{}HOY. 3) Revisar cantidad de rutas y secuencia de paradas.}{HTTP 200 agrupando ambos pedidos en una única propuesta ordenada cronológicamente (10:\allowbreak{}00, luego 10:\allowbreak{}29).}{Separación de 29 min \ensuremath{\leq} 30 se consolida en una ruta sin exceder capacidad máxima de 4.}
\caso{CP-019}{\textbf{Agrupación por extremos con separación exacta de 30 min.} NR-18, NR-22; NP-08 (C13.3, Media)}{Token ADMIN; 2 pedidos PENDING elegibles hoy con horas 10:\allowbreak{}00 y 10:\allowbreak{}30 (-05:\allowbreak{}00), diferencia exacta 30 minutos.}{1) Verificar margen=\allowbreak{}30. 2) Ejecutar GET /\allowbreak{}api/\allowbreak{}v1/\allowbreak{}route-proposals?date=\allowbreak{}HOY. 3) Inspeccionar propuestas y timeWindow\allowbreak{}Margin\allowbreak{}Minutes retornado.}{HTTP 200 agrupando ambos pedidos en una sola ruta al cumplirse la condición inclusiva max - min \ensuremath{\leq} 30.}{Igualdad exacta al límite (30 == 30) incluida en el mismo grupo con orden cronológico.}
\caso{CP-020}{\textbf{Separación de rutas por extremos con 31 minutos.} NR-18, NR-22; NP-08 (C13.4, Media)}{Token ADMIN; 2 pedidos PENDING elegibles hoy con horas 10:\allowbreak{}00 y 10:\allowbreak{}31 (-05:\allowbreak{}00), diferencia 31 minutos.}{1) Verificar margen=\allowbreak{}30. 2) Ejecutar GET /\allowbreak{}api/\allowbreak{}v1/\allowbreak{}route-proposals?date=\allowbreak{}HOY. 3) Verificar separación de grupos en proposals.}{HTTP 200 devolviendo 2 rutas propuestas separadas (1 parada cada una) por exceder el margen de 30 minutos.}{Separación de 31 min (\ensuremath{>} 30) impide la fusión de pedidos en la misma propuesta.}
\caso{CP-021}{\textbf{Cambio de margen horario (5..180) y desempate OSRM.} NR-18, NR-22; NP-08 (C13.5, Media)}{Token ADMIN; PATCH admin/\allowbreak{}shop-settings con routePlanning.\allowbreak{}timeWindow\allowbreak{}Margin\allowbreak{}Minutes=\allowbreak{}60; casos límite inválidos 4 y 181.}{1) Probar márgenes 4 y 181 \ensuremath{\rightarrow} 400. 2) Actualizar margen a 60 vía PATCH /\allowbreak{}api/\allowbreak{}v1/\allowbreak{}admin/\allowbreak{}shop-settings. 3) Consultar GET /\allowbreak{}api/\allowbreak{}v1/\allowbreak{}route-proposals?date=\allowbreak{}HOY.}{4 y 181 rechazados (400); con 60 min los 3 pedidos (10:\allowbreak{}00-10:\allowbreak{}50) se agrupan y aplican desempate OSRM durationsS primario / distancesM secundario.}{Margen validado en rango 5..180, persistido en settings y reflejado en nuevas propuestas con métricas viales coherentes.}
\caso{CP-022}{\textbf{Creación de pedido sin fecha programada (scheduledDeliveryDate nulo).} NR-10; NP-15 (C26.1, Media)}{Token CLIENT-A; payload POST /\allowbreak{}api/\allowbreak{}v1/\allowbreak{}orders con deliveryAddress, items y coordenadas válidas, omitiendo scheduledDeliveryDate.}{1) Enviar POST /\allowbreak{}api/\allowbreak{}v1/\allowbreak{}orders sin scheduledDeliveryDate. 2) Consultar GET /\allowbreak{}api/\allowbreak{}v1/\allowbreak{}orders/\allowbreak{}:\allowbreak{}id.}{HTTP 201; pedido creado con status=\allowbreak{}PENDING y scheduledDeliveryDate=\allowbreak{}null sin asignar fecha automática.}{Contrato opcional preservado: scheduledDeliveryDate permanece null y stock queda reservado.}
\caso{CP-023}{\textbf{Exclusión en propuestas y 409 en asignación sin fecha.} NR-10; NP-15 (C26.2, Media)}{Token ADMIN; pedido PENDING con scheduledDeliveryDate=\allowbreak{}null; fecha operativa HOY; driverId disponible.}{1) Identificar el pedido sin fecha en el listado administrativo. 2) Consultar propuestas de HOY. 3) Intentar confirmar y asignar vía POST /\allowbreak{}routes y POST /\allowbreak{}routes/\allowbreak{}assign.}{El pedido se identifica con fecha nula, no figura en propuestas y ambos intentos de asignación responden 409.}{Pedido identificable, no propuesto ni asignado; sin nueva ruta ni OrderClaim.}
\caso{CP-024}{\textbf{Idempotencia Shop checkout con replay 201.} NR-01, NR-05, NR-10 (C30.1, Alta)}{Token CLIENT-A; quoteToken vigente; Idempotency-Key UUIDv4 fijo y payload idéntico en dos solicitudes consecutivas.}{1) Enviar POST /\allowbreak{}api/\allowbreak{}v1/\allowbreak{}shop/\allowbreak{}checkout con Idempotency-Key K1. 2) Reenviar exactamente el mismo POST con K1 y mismo payload. 3) Consultar pedidos, recibos, stock y saldo.}{Ambas llamadas retornan HTTP 201 con idéntico orderId y receipt.\allowbreak{}id (replay idempotente).}{Solo existe un pedido, un comprobante, un descuento de stock y un débito de saldo demo.}
\caso{CP-025}{\textbf{Recuperación Shop ante timeout con mismo Idempotency-Key.} NR-01, NR-05, NR-10 (C30.2, Alta)}{Token CLIENT-A; quoteToken vigente; Idempotency-Key UUIDv4 K2 con payload válido bajo corte simulado de respuesta HTTP.}{1) Enviar POST /\allowbreak{}api/\allowbreak{}v1/\allowbreak{}shop/\allowbreak{}checkout con K2 provocando corte/\allowbreak{}timeout tras commit en servidor. 2) Reintentar con mismo K2 y payload. 3) Verificar stock y saldo.}{El reintento retorna HTTP 201 con el resultado original sin duplicar la transacción.}{Recuperación consistente: 1 pedido, 1 comprobante, 1 descuento de stock y 1 débito tras corte de red.}
\caso{CP-026}{\textbf{Conflicto Idempotency-Key con payload distinto (409).} NR-01, NR-05, NR-10 (C30.3, Alta)}{Token CLIENT-A; mismo Idempotency-Key UUIDv4 K3 reutilizado con payload alterado (cantidad o dirección distinta).}{1) Registrar conteo de pedidos, recibos, stock y saldo tras K3 original. 2) Enviar POST /\allowbreak{}api/\allowbreak{}v1/\allowbreak{}shop/\allowbreak{}checkout con K3 y payload distinto. 3) Reconsultar estado.}{HTTP 409 con error.\allowbreak{}code=\allowbreak{}IDEMPOTENCY\_KEY\_REUSED rechazando la reutilización conflictiva.}{No se crea segunda venta ni comprobante, ni se altera stock o saldo respecto al estado previo.}
\caso{CP-027}{\textbf{Liberación auditada de DRAFT caducado (RELEASE\_DRAFT).} NR-16, NR-20; NP-07 (C11.5, Alta)}{Token ADMIN; ruta DRAFT de ayer con pedido de ayer; driverId disponible; requestId UUIDv4 y reason válido.}{1) Intentar asignar la ruta y ejecutar RELEASE\_DRAFT. 2) Verificar fecha original, auditoría y replay. 3) Registrar acuerdo mediante CONTACT, aplicar RESCHEDULE con fecha futura y consultar propuestas de esa fecha.}{Asignación responde 409 ROUTE\_DATE\_NOT\_TODAY; se liberan paradas y claim sin cambiar la fecha. Tras acuerdo y reprogramación, el pedido aparece en propuestas de la nueva fecha.}{Pedido PENDING sin claim activo; no se reprograma automáticamente; auditoría y replay idempotente; elegibilidad posterior a la reprogramación.}
\end{disenio}
